\section{List of Findings}
\label{ListOfFindings}
\pagestyle{mypagestyle}
\begin{spacing}{1.75}
	
	\MUFGfont{
		This section lists findings identified during the independent review of the FX simulation model \footnote{(Model document 1.2, Section 2.4)}
		and its calibration approach \footnote{(Model document 1.1, Section 3.3).}
	}
	
	\subsection{List of New Findings}
	
	\begin{finding}
		\begin{table}[!h]
			\arrayrulecolor{mufggrey}
			\caption{\MUFGfont{Finding 01: Constant-vol GBM may under-represent tail risk and regime shifts for PFE}}
			\vspace{5pt}
			\begin{tabular}{| p{4.06cm}|p{4.06cm}|p{4.06cm}|p{4.06cm}|}
				\hline
				\multicolumn{2}{|l|}{\cellcolor{mufgred}\color{white}{\textbf{Finding}}} &
				\multicolumn{2}{l|}{\cellcolor{mufgred}\color{white}{\textbf{Recommendation}}} \\ \hline
				\multicolumn{2}{|p{8.12cm}|}{\MUFGfont{
						FX spots are simulated via GBM with time-homogeneous volatility and lognormal distribution.
						This specification does not capture volatility clustering, fat tails, or smile dynamics and may understate tail exposures during stress regimes,
						particularly when volatility is calibrated from a fixed historical window.
				}} &
				\multicolumn{2}{p{8.12cm}|}{\MUFGfont{
						Introduce mandatory sensitivity tests (vol window, return frequency, stressed scaling) and define a monitoring framework
						(e.g., volatility breaks, exceedance counts, PIT tests) with escalation triggers. Document explicit model limitations and usage guidance.
				}} \\ \hline
				\MUFGfont{\textbf{Risk Rating}} & \MUFGfont{Medium} &
				\MUFGfont{\textbf{Owner}: Model Owner} & \MUFGfont{\textbf{Date}: TBD}\\ \hline
			\end{tabular}
		\end{table}
	\end{finding}
	
	
	 
	 
	 \begin{finding}
	 	\begin{table}[!h]
	 		\arrayrulecolor{mufggrey}
	 		\caption{\MUFGfont{Finding 01 (Expanded): Constant-vol GBM may under-represent tail risk and regime shifts for PFE}}
	 		\vspace{5pt}
	 		\begin{tabular}{| p{4.06cm}|p{4.06cm}|p{4.06cm}|p{4.06cm}|}
	 			\hline
	 			\multicolumn{2}{|l|}{\cellcolor{mufgred}\color{white}{\textbf{Finding}}} &
	 			\multicolumn{2}{l|}{\cellcolor{mufgred}\color{white}{\textbf{Recommendation}}} \\ \hline
	 			
	 			\multicolumn{2}{|p{8.12cm}|}{\MUFGfont{
	 					\textbf{Model design.} FX spots are simulated as GBM with time-homogeneous volatility, implying normally distributed log-returns with independent increments. \par
	 					\textbf{Limitation.} This specification cannot reproduce empirically observed volatility clustering, regime shifts, excess kurtosis (fat tails), or risk-neutral skew/smile effects. \par
	 					\textbf{PFE impact.} During stressed regimes (or when volatility is calibrated from a benign historical window), GBM may compress extreme FX moves and therefore understate high-quantile exposure metrics (e.g., PFE\textsubscript{99}) and stressed EPE, especially for FX-sensitive instruments (e.g., XCCY swaps, non-USD collateral, FX conversions) and for portfolios where correlation-driven tail co-moves matter.
	 					
	 					
	 			}} &
	 			
	 			\multicolumn{2}{p{8.12cm}|}{\MUFGfont{
	 					
	 					\textbf{Mandatory sensitivity tests:} (i) volatility window length (1y/3y/5y), (ii) return frequency (1d vs 5d), (iii) stressed scaling overlays (e.g., max(current, stress-percentile)). \par
	 					\textbf{Tail validation:} rolling exception tests at 95\% 99\% for FX returns and/or exposure proxies; exceedance counts with escalation thresholds. \par
	 					\textbf{Monitoring:} volatility break detection (control charts / break tests), PIT/standardized residual tests, and regime indicators. \par
	 					\textbf{Governance:} document explicit limitations and define portfolio-level materiality triggers requiring overlays or enhanced modelling.
	 					
	 			}} \\ \hline
	 			
	 			\MUFGfont{\textbf{Risk Rating}} & \MUFGfont{Medium (upgrade to High if materiality triggers hold)} &
	 			\MUFGfont{\textbf{Owner}: Model Owner} & \MUFGfont{\textbf{Date}: TBD}\\ \hline
	 		\end{tabular}
	 	\end{table}
	 \end{finding}
	 
	 
	 
	 \begin{finding}
	 	\begin{table}[!h]
	 		\arrayrulecolor{mufggrey}
	 		\caption{\MUFGfont{Finding 02 (Expanded): Volatility calibration from 5-day returns requires justification, MPOR linkage, and stability evidence}}
	 		\vspace{5pt}
	 		\begin{tabular}{| p{4.06cm}|p{4.06cm}|p{4.06cm}|p{4.06cm}|}
	 			\hline
	 			\multicolumn{2}{|l|}{\cellcolor{mufgred}\color{white}{\textbf{Finding}}} &
	 			\multicolumn{2}{l|}{\cellcolor{mufgred}\color{white}{\textbf{Recommendation}}} \\ \hline
	 			
	 			\multicolumn{2}{|p{8.12cm}|}{\MUFGfont{
	 					\textbf{Calibration design.} The model calibrates $\sigma_k$ from a 3-year history of 5-day log returns, ignoring drift. \par
	 					\textbf{Limitation.} The 5-day choice embeds a specific horizon and implicitly assumes $\sigma_{1d}=\sigma_{5d}/\sqrt{5}$ for daily stepping. Under volatility clustering, this mapping may be biased. \par
	 					\textbf{PFE impact.} The approach may smooth short-term volatility spikes that can be material for MPOR-driven exposure, and may become stale or non-conservative during regime shifts.
	 			}} &
	 			
	 			\multicolumn{2}{p{8.12cm}|}{\MUFGfont{
	 					\textbf{Justification:} provide rationale linking 5-day returns to MPOR, liquidity, and operational stability. \par
	 					\textbf{Stability evidence:} compare $\sigma$ across (1d vs 5d returns), window lengths (1y/3y/5y), and EWMA alternatives; quantify impact on PFE metrics for representative portfolios. \par
	 					\textbf{Controls:} introduce volatility break triggers, minimum conservative floors, and define override governance (approval, duration, audit trail). \par
	 					\textbf{Acceptance criteria:} define tolerances for $\sigma$ drift, exception rates, and PFE sensitivity limits.
	 			}} \\ \hline
	 			
	 			\MUFGfont{\textbf{Risk Rating}} & \MUFGfont{Medium} &
	 			\MUFGfont{\textbf{Owner}: Model Owner} & \MUFGfont{\textbf{Date}: TBD}\\ \hline
	 		\end{tabular}
	 	\end{table}
	 \end{finding}
	 
	 
	 
	 %=========================================================
	 % Finding 02 — Mathematical rationale for 5d calibration
	 %=========================================================
	 \subsubsection{Horizon consistency of volatility estimation (1d vs 5d) \label{sec:fx_vol_scaling}}
	 
	 Let $S_t$ denote the FX spot. Define log-returns over horizon $\Delta$:
	 \begin{equation}
	 	r_{t,\Delta} \;=\; \ln\!\left(\frac{S_{t+\Delta}}{S_t}\right).
	 \end{equation}
	 
	 \paragraph{Daily (1d) volatility estimator.}
	 From a window of $N$ daily observations, the sample variance estimator is
	 \begin{equation}
	 	\widehat{\sigma}^2_{1d}
	 	=
	 	\frac{1}{N-1}\sum_{i=1}^{N}\left(r_{t_i,1d}-\overline{r}_{1d}\right)^2,
	 	\qquad
	 	\overline{r}_{1d}=\frac{1}{N}\sum_{i=1}^{N} r_{t_i,1d}.
	 \end{equation}
	 For short horizons, drift is negligible relative to diffusion and $\overline{r}_{1d}\approx 0$ is commonly assumed in PFE contexts.
	 
	 \paragraph{5-day (5d) volatility estimator.}
	 From $M$ observations of 5-day log-returns, either overlapping
	 \begin{equation}
	 	r_{t,5d}=\ln\!\left(\frac{S_{t}}{S_{t-5}}\right),
	 \end{equation}
	 or non-overlapping (every 5th day), the corresponding sample variance is
	 \begin{equation}
	 	\widehat{\sigma}^2_{5d}
	 	=
	 	\frac{1}{M-1}\sum_{j=1}^{M}\left(r_{t_j,5d}-\overline{r}_{5d}\right)^2.
	 \end{equation}
	 This estimator targets variability over a 5-day horizon.
	 
	 \paragraph{Square-root-of-time mapping (model assumption).}
	 Under i.i.d.\ increments (e.g. constant-vol GBM), the 5-day return is the sum of five independent 1-day returns:
	 \begin{equation}
	 	r_{t,5d}=\sum_{k=0}^{4} r_{t-k,1d},
	 	\qquad
	 	\mathrm{Var}(r_{t,5d}) = \sum_{k=0}^{4}\mathrm{Var}(r_{t-k,1d}) = 5\,\sigma^2_{1d},
	 \end{equation}
	 so that
	 \begin{equation}
	 	\sigma_{5d}=\sqrt{5}\,\sigma_{1d},
	 	\qquad
	 	\sigma_{1d}=\frac{\sigma_{5d}}{\sqrt{5}}.
	 \end{equation}
	 When 5-day calibration is used but the Monte Carlo engine steps daily, the implementation implicitly relies on this scaling identity.
	 
	 \paragraph{Implied daily volatility from 5d calibration.}
	 Define the implied 1-day volatility:
	 \begin{equation}
	 	\widehat{\sigma}^{(\mathrm{implied})}_{1d,5d}
	 	=
	 	\frac{\widehat{\sigma}_{5d}}{\sqrt{5}}.
	 \end{equation}
	 
	 \paragraph{Horizon-consistency metric.}
	 Define the relative discrepancy
	 \begin{equation}
	 	\Delta_{\sigma}
	 	=
	 	\frac{\widehat{\sigma}^{(\mathrm{implied})}_{1d,5d}-\widehat{\sigma}_{1d}}{\widehat{\sigma}_{1d}}.
	 \end{equation}
	 Under i.i.d.\ increments and correct scaling, $\Delta_{\sigma}$ should be close to zero (up to estimation noise). Persistent deviations (especially negative $\Delta_{\sigma}$) indicate that the 5-day choice smooths short-horizon volatility and may lag regime shifts, which can be material for MPOR-driven tail exposure metrics (e.g.\ PFE$_{99}$).
	 
	 \paragraph{Note on overlapping 5d returns.}
	 If 5-day returns are computed daily (overlapping), successive $r_{t,5d}$ share four out of five daily moves, inducing serial dependence in the calibration sample. This affects effective sample size and the responsiveness of $\widehat{\sigma}_{5d}$ to breaks; the chosen convention should therefore be explicitly documented and justified.
	 
	 
	
	
	\begin{finding}
		\begin{table}[!h]
			\arrayrulecolor{mufggrey}
			\caption{\MUFGfont{Finding 03: Measure consistency and drift definition should be explicitly stated and tested}}
			\vspace{5pt}
			\begin{tabular}{| p{4.06cm}|p{4.06cm}|p{4.06cm}|p{4.06cm}|}
				\hline
				\multicolumn{2}{|l|}{\cellcolor{mufgred}\color{white}{\textbf{Finding}}} &
				\multicolumn{2}{l|}{\cellcolor{mufgred}\color{white}{\textbf{Recommendation}}} \\ \hline
				\multicolumn{2}{|p{8.12cm}|}{\MUFGfont{
						Section 2.4 sets the mean function $g_k^{FX}(t)$ from discount factors, implying forward-implied drift.
						The documentation should explicitly link the chosen drift to the simulation measure used in PFE and verify consistency with numeraire/discounting
						in downstream pricing/valuation under scenarios.
				}} &
				\multicolumn{2}{p{8.12cm}|}{\MUFGfont{
						Add an explicit statement of the measure/numeraire used for simulation and valuation, and implement unit tests:
						(1) forward-martingale checks (discounted FX expectations),
						(2) cross-rate consistency checks derived from USD legs.
				}} \\ \hline
				\MUFGfont{\textbf{Risk Rating}} & \MUFGfont{Low/Medium} &
				\MUFGfont{\textbf{Owner}: Model Owner} & \MUFGfont{\textbf{Date}: TBD}\\ \hline
			\end{tabular}
		\end{table}
	\end{finding}
	
	
	\newpage 
	
	\section{Finding 01 (Expanded): Constant-Volatility GBM and Tail Risk in FX PFE Simulation}
	\label{Finding01Expanded}
	\pagestyle{mypagestyle}
	\begin{spacing}{1.75}
		
		\subsection{Context and Scope}
		This finding relates to the modelling choice used for FX spot dynamics within the PFE simulation engine, as described in Model document~1.2, Section~2.4 (FX simulation), and its associated calibration approach described in Model document~1.1, Section~3.3.
		
		The model simulates FX spot rates (all currencies versus USD) using a geometric Brownian motion (GBM) with time-homogeneous volatility, where the drift is implied by the FX forward curve constructed from domestic and foreign discount factors. Volatility parameters are calibrated from a fixed historical window of 5-day log returns.
		
		This finding assesses whether this specification is adequate for capturing tail exposure risk under PFE, particularly during stressed market regimes.
		
		\subsection{Model Specification}
		
		\subsubsection{FX Dynamics under the Model}
		Let $S_t$ denote the FX spot rate (USD per unit of foreign currency). Under the model, $S_t$ follows a GBM of the form:
		\begin{equation}
			\frac{\mathrm{d}S_t}{S_t} = \mu(t),\mathrm{d}t + \sigma,\mathrm{d}W_t,
		\end{equation}
		where $\sigma$ is a constant volatility parameter and $W_t$ is a standard Brownian motion.
		
		The drift $\mu(t)$ is not estimated statistically but is implied by the FX forward curve through the deterministic mean function:
		\begin{equation}
			g(t) = S_0\frac{DF_{\text{FOR}}(0,t)}{DF_{\text{USD}}(0,t)},
		\end{equation}
		so that the simulated FX rate is expressed as
		\begin{equation}
			S_t = g(t)\exp\left(X_t - \tfrac{1}{2}\sigma^2 t\right), \qquad X_t = \sigma W_t.
		\end{equation}
		
		\subsubsection{Distributional Implications}
		Under this specification, log-returns over a time step $\Delta t$ satisfy:
		\begin{equation}
			\Delta \ln S \sim \mathcal{N}\left((\mu-\tfrac{1}{2}\sigma^2)\Delta t,, \sigma^2\Delta t\right).
		\end{equation}
		As a result:
		\begin{itemize}
			\item log-returns are normally distributed (thin tails);
			\item variance scales linearly with time;
			\item increments are independent and identically distributed.
		\end{itemize}
		
		\subsection{Empirical Properties of FX Returns}
		Extensive empirical evidence shows that FX returns exhibit:
		\begin{itemize}
			\item \textbf{Volatility clustering}: conditional heteroskedasticity with persistent high- and low-volatility regimes;
			\item \textbf{Regime shifts}: abrupt changes in volatility levels during crises or macro events;
			\item \textbf{Excess kurtosis}: heavier tails than implied by a Gaussian distribution;
			\item \textbf{Skew and smile effects}: asymmetry and curvature in the risk-neutral distribution.
		\end{itemize}
		
		These features are well documented in the FX literature and are not captured by a constant-volatility GBM.
		
		\subsection{Implications for PFE and Exposure Tails}
		
		\subsubsection{Tail Behaviour under GBM}
		For a given horizon $T$, the $q$-quantile of $S_T$ under GBM is given by:
		\begin{equation}
			Q_q(S_T) = g(T)\exp\left(\sigma\sqrt{T},\Phi^{-1}(q) - \tfrac{1}{2}\sigma^2 T\right),
		\end{equation}
		where $\Phi^{-1}$ is the inverse standard normal distribution.
		
		This expression highlights that tail behaviour is entirely driven by $\sigma$ and the Gaussian quantile $\Phi^{-1}(q)$. Any underestimation of $\sigma$ or omission of fat-tail behaviour directly compresses extreme quantiles.
		
		\subsubsection{Why This Matters for PFE}
		Potential Future Exposure focuses explicitly on high quantiles of the exposure distribution (e.g., 95%, 99%). FX enters PFE through:
		\begin{itemize}
			\item cross-currency swaps and foreign currency legs;
			\item FX forwards and notional exchanges;
			\item collateral currency mismatches and FX discounting effects;
			\item correlation-driven tail co-movements with rates and credit.
		\end{itemize}
		
		Even when individual trades are linear in FX, portfolio-level exposures can become highly nonlinear in the tail. Under constant-volatility GBM, these tail scenarios may be systematically under-represented, particularly when volatility is calibrated from a benign historical window.
		
		\subsection{Interaction with Volatility Calibration}
		
		Volatility is calibrated from a three-year history of 5-day log returns. This introduces two compounding effects:
		\begin{enumerate}
			\item \textbf{Smoothing effect}: 5-day returns smooth short-term volatility spikes and may understate sudden FX moves relevant for MPOR-driven exposure;
			\item \textbf{Regime averaging}: a fixed historical window blends calm and stressed regimes, potentially producing a volatility estimate that is not conservative during stress.
		\end{enumerate}
		
		Under volatility clustering, the standard square-root-of-time scaling ($\sigma_{1d}=\sigma_{5d}/\sqrt{5}$) may not hold, further distorting tail risk.
		
		\subsection{Materiality Assessment}
		This finding is particularly material when one or more of the following conditions hold:
		\begin{itemize}
			\item large FX notionals or long-dated cross-currency products;
			\item collateral posted in a currency different from the exposure currency;
			\item significant FX--rates correlation within the portfolio;
			\item portfolios sensitive to tail FX moves rather than average volatility.
		\end{itemize}
		
		If such conditions are present, the risk rating of this finding should be escalated from \textit{Medium} to \textit{High}.
		
		\subsection{Validation Tests and Evidence Required}
		
		To support or mitigate this finding, the following analyses are required:
		\begin{itemize}
			\item \textbf{Return-level diagnostics}: empirical kurtosis, tail quantiles, and exceedance tests for 1-day and 5-day FX returns;
			\item \textbf{Stability analysis}: rolling volatility estimates across multiple window lengths and return frequencies;
			\item \textbf{Regime analysis}: comparison of calibrated volatility and implied tails across calm and stressed sub-periods;
			\item \textbf{PFE sensitivity}: impact of alternative volatility assumptions (1d, 5d, EWMA, stressed scaling) on portfolio-level PFE metrics.
		\end{itemize}
		
		\subsection{Conclusion of Finding}
		The use of a constant-volatility GBM for FX simulation is a standard and tractable choice for PFE engines. However, from a tail-risk perspective, it represents a structural simplification that can understate extreme FX moves and, consequently, high-quantile exposure measures.
		
		Absent strong empirical evidence and compensating controls, this limitation constitutes a material model risk. The finding can be mitigated through robust sensitivity testing, conservative volatility overlays, and ongoing monitoring, without necessarily replacing the core GBM framework.
		
	\end{spacing}
	\clearpage
	
	
\end{spacing}
\clearpage
